\documentclass[12pt, letterpaper]{article}
\usepackage[utf8]{inputenc}
\usepackage{geometry}
\geometry{margin=1in} % Mandatory 1-inch margins
\usepackage{titlesec}
\usepackage{enumitem}
\usepackage{fancyhdr} % Mandatory for Last Name + Page Number headers
\usepackage{setspace} % Mandatory for Double Spacing

% Mandatory Double Spacing
\doublespacing

% Mandatory Header setup for pages 2-5
\pagestyle{fancy}
\fancyhf{}
\renewcommand{\headrulewidth}{0pt}
\fancyhead[R]{Lee \thepage} % Last Name and Page Number

% Section formatting
\titleformat{\section}{\large\bfseries\uppercase}{}{0em}{}[\titlerule]
\titleformat{\subsection}{\bfseries}{}{0em}{}

\begin{document}

% Page 1 Header: Full Name and Project Title (Mandatory)
\thispagestyle{empty}
\begin{center}
    \textbf{JUNSEONG LEE} \\
    \textbf{THE RADIOMICS-MICROBIOME AXIS LAYER} \\
    Extending SanoMap via Agentic AI — Delivered System
\end{center}

\section{A statement of the research problem}

Biomedical knowledge graphs derive their clinical value from the completeness of the evidence they encode. Recent advances in natural language processing—primarily the emergence of transformer-based models and generative AI—have enabled large-scale construction of these graphs by systematically mining verified scientific literature. The SanoMap project's MINERVA platform (Microbiomes Network for Research and Visualization Atlas) demonstrates this capability well: it maps associations between gut microbes and systemic diseases at scale, producing a clinically navigable graph from thousands of curated biomedical papers.

Yet MINERVA's graph encodes a structurally incomplete picture. It connects microbes directly to diseases, but omits the quantitative, tissue-level intermediate phenotypes that explain the biological mechanism linking them. Radiomics and body composition imaging provide exactly these intermediaries. Features such as skeletal muscle index, visceral adipose tissue fraction, liver fat fraction, and GLCM texture signatures are extracted directly from CT and MRI scans and are well-established biomarkers for conditions like sarcopenia, nonalcoholic fatty liver disease, and colorectal cancer. Importantly, the scientific literature increasingly demonstrates that gut microbiome composition measurably influences these specific tissue-level phenotypes---making radiomics and body composition features natural intermediary nodes in any complete microbiome-disease knowledge graph.

A second structural gap concerns the evidence modality itself. Standard text-mining pipelines extract findings that authors highlight in prose, but much of the quantitative correlation data linking microbial taxa to imaging feature values is encoded in figures. Gradient correlation heatmaps, for instance, display matrices of r-values linking dozens of taxa to radiomic measurements—data that NLP alone cannot parse. This constitutes a systematic evidence loss in text-only pipelines.

This work addresses both gaps. The central research questions are: (1) Can a hybrid AI pipeline combining large language models and vision-language models, with deterministic verification, reliably extract quantitative correlation coefficients directly from medical figures? And (2) does adding radiomics as an ontological layer---filling the microbe-to-imaging-phenotype-to-disease path---meaningfully increase the clinical completeness of the MINERVA knowledge graph?

\section{The Radiomics Intermediary Layer: Concept and Validation}

\subsection{The Core Idea}

MINERVA's current graph encodes the relation \texttt{(Microbe)$\to$[ASSOCIATED\_WITH]$\to$(Disease)}. The delivered system argues that this edge is clinically underspecified. The biologically accurate path is:

\begin{center}
\texttt{Microbe} $\to$ \textit{correlates with intermediate phenotype} $\to$ \texttt{RadiomicFeature / BodyCompositionFeature} $\to$ \textit{predicts/associates with} $\to$ \texttt{Disease}
\end{center}

This three-node path is not a theoretical addition. Clinical studies document it concretely: \textit{Prevotella nigrescens} correlates with GLCM Correlation texture features on CT at $r = 0.95$ (PMC10605408); gut dysbiosis in cirrhosis is characterized by enrichment of Proteobacteria alongside measurable changes in liver surface nodularity and visceral fat; \textit{Bifidobacterium} depletion co-occurs with elevated visceral adipose tissue and downstream obesity. These are not inferred associations—they are quantitatively measured, directly evidenced relations that a text-only microbe-to-disease edge collapses into a single, mechanism-blind link.

The intermediary node layer proposed here makes this mechanism explicit and traversable in the graph, enabling queries that are not possible in MINERVA's current schema: ``Which radiomic features of the liver are associated with both Fusobacterium enrichment and hepatocellular carcinoma?'' or ``What CT texture signature mediates the association between gut dysbiosis and sarcopenia?''

\subsection{Graph Schema}

The delivered system extends MINERVA with the following node types:

\begin{itemize}
    \item \texttt{RadiomicFeature} — texture, morphological, and intensity imaging features (GLCM, wavelet, first-order, shape)
    \item \texttt{BodyCompositionFeature} — quantitative tissue phenotypes from imaging: sarcopenia, visceral adipose tissue, skeletal muscle index, liver fat fraction, myosteatosis
    \item \texttt{BodyLocation} — anatomical measurement sites (liver, colon, lumbar spine, brain, kidney, heart, and others)
    \item \texttt{ImagingModality} — acquisition modalities (CT, MRI, PET, DXA, ultrasound, X-ray)
    \item \texttt{ImageRef} — verified figure references with provenance from the Vision Track
    \item \texttt{Microbe} and \texttt{MicrobialSignature} — microbial entity nodes, consistent with MINERVA's upstream schema
    \item \texttt{Disease} — disease entity nodes from literature extraction
\end{itemize}

The direct-evidence edge policy, designed to be semantically compatible with MINERVA's edge types, is:

\begin{itemize}
    \item \texttt{(Microbe)-[:CORRELATES\_WITH]->(RadiomicFeature)} — quantitatively verified
    \item \texttt{(Microbe)-[:CORRELATES\_WITH]->(BodyCompositionFeature)} — quantitatively verified
    \item \texttt{(RadiomicFeature)-[:ASSOCIATED\_WITH]->(Disease)} — text co-mention, associative evidence
    \item \texttt{(BodyCompositionFeature)-[:ASSOCIATED\_WITH]->(Disease)} — text co-mention, associative evidence
    \item \texttt{(RadiomicFeature)-[:MEASURED\_AT]->(BodyLocation)} — anatomical grounding
    \item \texttt{(RadiomicFeature)-[:ACQUIRED\_VIA]->(ImagingModality)} — modality provenance
    \item \texttt{(ImagingModality)-[:REPRESENTED\_BY]->(ImageRef)} — figure-level citation
    \item \texttt{(Microbe)-[:POSITIVELY\_CORRELATED\_WITH / :NEGATIVELY\_CORRELATED\_WITH]->(Disease)} — signed microbe-disease relation, compatible with MINERVA's upstream edges
\end{itemize}

The distinction between \texttt{CORRELATES\_WITH} and \texttt{ASSOCIATED\_WITH} is deliberate. Quantitatively verified edges (from the Vision Track, where an r-value is deterministically confirmed against a pixel-level color scale) use \texttt{CORRELATES\_WITH}. Text-derived edges, which reflect co-mention associations without quantitative grounding, use \texttt{ASSOCIATED\_WITH}. This preserves evidence-level semantics across the graph.

\subsection{Data Sourcing}

Data was sourced from PubMed Central Open Access using a multi-lane query strategy targeting four clinical domains where the microbiome-radiomics connection is most evidenced:

\begin{enumerate}
    \item \textbf{Strict radiomics + microbiome}: papers explicitly co-mentioning radiomic features (GLCM, wavelet, first-order) with microbiome markers
    \item \textbf{Body composition + microbiome}: CT/MRI body composition papers (sarcopenia, visceral adipose tissue, skeletal muscle index) with gut microbiome co-mention
    \item \textbf{Liver radiomics + microbiome}: hepatic steatosis, NASH/NAFLD, liver fat fraction, liver surface nodularity with microbiota
    \item \textbf{Bone density + microbiome}: DXA-measured BMD, osteoporosis, trabecular bone score with gut microbiome
    \item \textbf{Lung CT phenotypes + microbiome}: emphysema, COPD CT metrics, pulmonary fibrosis with microbiota
    \item \textbf{Colorectal imaging + microbiome}: CT radiomics of colorectal tumors, tumor texture, with gut/intratumoral microbiome
\end{enumerate}

The full corpus spans \textbf{1,016 papers} (640 from the initial three lanes plus 376 net-new from lanes 4--6), with open-access full text retrieved for papers with PMCIDs.

\subsection{Two-Track Pipeline Architecture and the Five-Gate Acceptance Model}

The system operates on two evidence tracks that populate complementary parts of the graph; each track is governed by a stack of independent acceptance gates so that a candidate edge reaches the graph only when every applicable gate passes.

The \textbf{Text Track} extracts phenotype-to-disease associations at scale. A vocabulary-matched extraction stage maps diverse imaging phenotype terminology to canonical names. Microbial entity recognition uses the \texttt{d4data/biomedical-ner-all} model on Apple Silicon MPS; disease recognition uses the \texttt{en\_ner\_bc5cdr\_md} spaCy pipeline (BC5CDR). A shared span cleanup layer rejects wordpiece artifacts, generic entity placeholders, and clause-prefix fragments before relation classification. Relations are extracted using Gemini 2.5 Flash-Lite via a self-consistency protocol: 7 samples per entity pair, accepted only at full agreement.

The \textbf{Vision Track} addresses the structured figure data gap. A vision-language model (Gemini 2.5 Flash-Lite) proposes r-values for microbial-entity-to-radiomic-feature correlations extracted from heatmap figures. Each proposal is cross-checked by a deterministic pixel-level verifier: the predicted value's color is checked for geometric consistency with the figure's gradient color scale legend, and acceptance requires a matching pixel support fraction above a calibrated threshold.

\textbf{Five-Gate Acceptance Model.} Tracks alone do not establish research-defensibility; structurally distinct gates do. Each candidate edge is evaluated against the following stack:
\begin{enumerate}
    \item \textbf{Retrieval (text)} --- replaces the original hand-curated 25-alias \texttt{\_FEATURE\_VOCAB} substring filter with a dense-retrieval index built from \texttt{thomas-sounack/BioClinical-ModernBERT-base} embeddings (mean-pooled with attention mask, L2-normalized, FAISS index where available, NumPy cosine fallback at the current corpus size of $\sim$30k sentences). Per-feature similarity threshold $\tau$ is calibrated per concept against labeled data. This gate addresses the recall ceiling of substring matching against paraphrastic feature mentions (``loss of muscle mass at L3'', ``diminished cross-sectional muscle area'').
    \item \textbf{Entity sanitization (UMLS TUI grounding)} --- microbial entity surface forms must ground to a microbe-class UMLS Semantic Type at similarity $\geq 0.85$. Accept set is \{T005 Virus, T007 Bacterium, T194 Archaeon, T204 Eukaryote\} with an explicit deny-list of non-microbe model-organism CUIs (\textit{Homo sapiens}, \textit{Mus musculus}, etc.). This gate rejects entity-type errors that the upstream NER produced but the substring filter and self-consistency gate missed --- specifically the published Edge \#5 case (\texttt{gut bacterial clpb-like gene function $\to$ body\_fat}), where a KEGG gene-function noun phrase was tagged as an organism. A live audit of the 8 accepted text-derived microbe$\to$feature edges drops 25\% (2/8): Edge \#5 at similarity 0.799 < 0.850, plus the typo \texttt{bacteriodetes} (no UMLS match). The 25\% drop rate sits below the 30\% recalibration threshold; the similarity floor remains 0.85.
    \item \textbf{Relation acceptance (text)} --- Gemini 2.5 Flash-Lite, 7-sample temperature-varied self-consistency, full-agreement required. The measured full-consistency rate on the expanded corpus is 0.916.
    \item \textbf{Verification (vision) --- dual verifier} --- single-modality pixel verification is structurally vulnerable to VLM monoculture: a same-family proposer + verifier shares correlated errors that self-consistency cannot detect. The Vision Track therefore runs an AND-consensus dual gate: the existing pixel HSV verifier plus an independent VLM verifier called with a verifier-only prompt at temperature 0. The default deployment uses local \texttt{qwen2.5vl:3b} (Q4\_K\_M, $\sim$3.2~GB) served via Ollama at \texttt{localhost:11434} --- chosen so that the verifier costs zero per inference and runs entirely within the 8~GB unified-memory ceiling on the development machine. Convergent validity is claimed at the modality level (pixel-data vs.\ language-level interpretation), not at the data level (both verifiers depend on image quality), and is documented as such in the limitations. Disagreement (XOR-pass) is preserved by routing to a review queue rather than silently rejected. A live smoke test on the two locally available qualifying figures returned AND-consensus PASS for both: \textit{Peptostreptococcus} $\leftrightarrow$ DLCO/VA\%pred at $r=-0.407$ (pixel distance 0.019; Qwen colour band ``light blue''); \textit{Haemophilus} $\leftrightarrow$ 4th~Ai at $r=-0.6$ (pixel distance 0.0; Qwen colour band ``blue'').
    \item \textbf{Evaluation (gold benchmark)} --- a 5-stratum sampler (\texttt{accepted\_edge}, \texttt{gemini\_rejected}, \texttt{vocab\_excluded}, \texttt{recall\_probe}, \texttt{random\_co\_occurrence}) drawn deterministically from the production candidate space produces a 66-row gold set. The annotation schema (v1.0) defines a 6-class primary label, three secondary labels (\texttt{evidence\_type} / \texttt{quantitative} / \texttt{confidence}), eight numbered edge-case decisions, and a 14-day intra-annotator temporal-relabeling protocol with Cohen's $\kappa$ acceptance threshold $\geq 0.80$ on the binary collapse. Pass-1 labels were generated under a computer-aided manual annotation workflow (LLM proposer, single human annotator review) with full per-row rationale captured for review traceability; Pass-2 is the human annotator's independent re-label after the temporal window, and Cohen's $\kappa$ on the two human passes (Pass-1 reviewed and overridden, Pass-2 independent) is the schema-clarity metric of record. Numbers in this section that rely on Pass-1 + Pass-2 are pending the temporal window's close.
\end{enumerate}

This gating model converts the original two-track pipeline from a proof-of-concept end-to-end run into a measurable artifact. Each gate is independently auditable; an edge enters the graph only when it survives every gate that applies to its evidence track.

\subsection{Validation: Evidence That the Idea Produces Meaningful Results}

Across the full 1,016-paper corpus, the pipeline extracted:

\begin{itemize}
    \item \textbf{5,721} phenotype text mentions across 30 clinically meaningful disease targets (from the initial 640-paper corpus; re-extraction from the expanded corpus is in progress)
    \item \textbf{74} \texttt{ASSOCIATED\_WITH} edges (phenotype-to-disease), validated by Gemini self-consistency (Gemini full-consistency rate: 0.916 on the expanded corpus)
    \item \textbf{29} signed microbe-disease pairs (\texttt{POSITIVELY\_CORRELATED\_WITH} / \texttt{NEGATIVELY\_CORRELATED\_WITH}): 14 positive, 15 negative, covering cirrhosis, NAFLD, colorectal cancer, inflammatory bowel disease, celiac disease, chronic HCV infection, and systemic sclerosis-associated pulmonary fibrosis
    \item \textbf{8} \texttt{CORRELATES\_WITH} edges (microbe-to-feature) after the 2026-05-07 retroactive vision-track audit (see § ``Vision Track Audit and Honest Scoping''): 1 from Vision Track (\textit{Prevotella nigrescens} $\to$ GLCM\_Correlation, $r=0.95$, PMC10605408\_g004 panel A, real Spearman correlation heatmap with $-1$ to $+1$ colorbar; pixel verifier distance 0.05, support fraction 1.0) and 7 from Text Track (Gemini self-consistency classification of microbe--feature co-occurrence sentences, 7/7 full agreement required). One pre-audit Vision Track edge (\textit{Firmicutes} $\to$ Total fat \%, claimed $r=-0.95$ at PMC6178902\_g0006) was retracted: direct image inspection found the cell is deep red (positive, magnitude $\approx +1.0$), the colorbar runs $-1.5$ to $+1.5$ (LFC scale, not Pearson $r$), and the original ``pixel-verified'' pass was a self-consistent fabrication on the proposer's hallucinated bounding box. Three of the surviving \texttt{CORRELATES\_WITH} edges close end-to-end three-hop paths through features that also carry \texttt{ASSOCIATED\_WITH} disease edges: \textit{Ruminococcus} $\to$ sarcopenia (37 diseases), \textit{Peptostreptococcus stomatis} $\to$ skeletal\_muscle\_index (7 diseases), and \textit{Eubacterium} $\to$ visceral\_adipose\_tissue (18 diseases) --- yielding \textbf{62 traversable microbe $\to$ feature $\to$ disease paths}
    \item \textbf{18} \texttt{BodyLocation} nodes, \textbf{5} \texttt{ImagingModality} nodes
    \item \textbf{191} total Neo4j export rows; \textbf{280} automated test checks passing across the full suite (was 156 prior to the structural-upgrade work; +105 new tests cover the UMLS validator, dense retrieval, dual verifier, gold-set sampler, and benchmark evaluator with no regressions on prior tests)
\end{itemize}

Two findings from the expanded corpus merit specific note. First, \textit{Faecalibacterium} depletion emerged as a convergent signal across six distinct disease associations in the new liver and colorectal lanes: chronic liver disease, colorectal cancer, inflammatory bowel disease, celiac disease, irritable bowel syndrome, and systemic sclerosis-associated pulmonary fibrosis. This cross-disease depletion pattern is consistent with Faecalibacterium's established role as a butyrate-producing protective commensal. Second, \textit{Fusobacterium nucleatum} enrichment in colorectal cancer was independently confirmed across four PMIDs from the colorectal imaging lane, validating the lane's signal quality. Both findings represent disease domains (colorectal oncology, liver fibrosis) where the microbiome-imaging-phenotype link is clinically documented---precisely the target of the radiomics intermediary layer.

The full Vision Track corpus run processed 74 PMCIDs (120-paper input, 74 with accessible figures), yielding 49 figures classified as qualifying topologies (heatmap, forest plot, scatter plot, dot plot). A two-stage relevance gate—checking whether each figure shows a direct quantitative association between a microbial entity and a radiomics or imaging feature—passed 13 of 49. Of these, 12 produced successful VLM proposals (1 returned a transient API error); 1 was verified by the pixel-level color scale validator. The verifier's specificity is consistent: non-heatmap figures (forest plots, scatter plots) pass the relevance gate but correctly fail the pixel-level verifier because they lack a gradient color scale, resulting in support fraction = 0. This is expected—the verifier targets gradient correlation matrices, and the pipeline correctly gates on that evidence standard.

\subsection{Vision Track Audit and Honest Scoping}

After the four-task structural upgrade landed, a retroactive audit of all Vision Track inputs (13 current proposals plus 1 historical graph edge) revealed that the dual-verifier consensus alone is structurally insufficient: the pixel HSV verifier and the independent VLM verifier both consume the proposer's bounding box and proposed $r$-value, so a self-consistent fabrication (proposer hallucinates both fields such that the predicted colour happens to match the colorbar at that location on the assumed default $\pm 1.0$ scale) passes AND-consensus silently. The pixel verifier additionally defaults its \texttt{observed\_range} to $[-1.0, +1.0]$ rather than reading the figure's tick labels, so heatmaps using log-fold-change ($\pm 1.5$) or other non-$r$ scales are validated as if they were Pearson $r$ matrices.

To close this hole, three deterministic pre-verifier gates were implemented in \texttt{src/vision\_gates.py} and run before AND-consensus is even attempted:
\begin{enumerate}
    \item \textbf{Caption gate} --- the figure caption must contain explicit Pearson / Spearman / correlation-coefficient vocabulary, and must \emph{not} contain auto-fail tokens (\texttt{log fold change}, \texttt{LFC}, \texttt{log2}, \texttt{z-score}, \texttt{differential abundance}). This deterministically rejects scatter plots, network diagrams, log-fold-change heatmaps, and other figures the proposer's topology classifier mis-tagged as heatmaps.
    \item \textbf{Colorbar detection gate} --- a gradient colorbar legend must be detectable in the image (reuses the existing legend detector from \texttt{verify\_heatmap}). Network and pathway diagrams without colorbars fail here.
    \item \textbf{Range-sanity gate} --- $|proposed\,r| \leq \max(|c_{\min}|, |c_{\max}|) + 0.05$. The colorbar bounds $(c_{\min}, c_{\max})$ are extracted by an additional focused VLM call (\texttt{extract\_colorbar\_range\_via\_vlm}) that reads the figure's actual tick labels rather than relying on the default $\pm 1.0$. Out-of-range hallucinations on tighter colorbars (e.g., a proposer claiming $r = +0.78$ on a $\pm 0.23$ colorbar) are caught here.
\end{enumerate}

On the audit batch ($n = 14$) this gate chain rejected 6 figures (3 caption fails, 2 range-sanity fails, 1 colorbar-detect fail), and the dual verifier subsequently flagged 3 of the remaining 8 for human review. The single historical Vision Track edge that survived (the \textit{Prevotella nigrescens} $\to$ GLCM\_Correlation reading on PMC10605408\_g004) was retained on the basis of (a) caption-confirmed Spearman correlation heatmap, (b) extracted $-1$ to $+1$ colorbar matching the proposed $r = 0.95$, (c) Session-4 pixel distance of 0.05 with support fraction 1.0, and (d) direct image inspection. The other historical Vision Track edge (\textit{Firmicutes} $\to$ Total fat \%, $r = -0.95$ at PMC6178902\_g0006) was retracted from the graph by \texttt{scripts/drop\_failed\_vision\_edges.py} after direct image inspection identified that the asterisked cell is in the figure's positive (red) hemisphere, the colorbar runs $-1.5$ to $+1.5$ (consistent with a log-fold-change scale rather than Pearson $r$), and the prior pass had been a circular consistency check.

This audit reframes the Vision Track contribution honestly: rather than the pre-audit ``9 \texttt{CORRELATES\_WITH} edges (2 vision + 7 text)'' framing, the post-audit graph carries 8 edges (1 vision + 7 text), and the methods section documents the dual-verifier's structural limitation alongside the pre-verifier gate chain that mitigates it. A residual failure mode that the gates do not catch is wrong-sign hallucination when the proposer's bounding box happens to point at a same-colour cell elsewhere in the figure; an additional sign-check gate (proposer-claimed sign versus VLM-observed colour hemisphere) is documented as the next-cheapest extension.

\section{Comparison with MINERVA and Iterative Gap Analysis}

The delivered system is built to extend MINERVA, not to replicate it from scratch. The following gaps exist relative to MINERVA's upstream methodology. Each is treated as an explicitly acknowledged design decision, not an error.

\subsection{Gap 1: Named Entity Recognition Models}

\textbf{MINERVA's approach:} A custom DistilBERT model (BNER2.0) fine-tuned on a proprietary corpus of 12,480 microbiome entity annotations across 23,174 sentences, achieving F1 = 0.914. The fine-tuned weights are not publicly released.

\textbf{This system's approach:} The \texttt{d4data/biomedical-ner-all} model (MACCROBAT-trained, 107 biomedical entity types) handles microbial NER; \texttt{en\_ner\_bc5cdr\_md} handles disease NER. A controlled 5-paper benchmark (2026-03-07) compared four substitute models; \texttt{d4data/biomedical-ner-all} produced the most stable relation output. The BENT-PubMedBERT-NER-Organism model, a publicly available organism-specific model, produced fewer usable pairs on the same subset.

\textbf{Residual risk:} The generic biomedical NER model likely has lower precision on gut microbiome taxonomy (phylum/genus/species disambiguation) than a corpus-specific fine-tune. This is partially compensated by the downstream span cleanup layer and the genus allowlist, which reject generic microbial terms and wordpiece fragments before relation extraction.

\textbf{Path to closure:} Two paths exist. First, if BNER2.0 weights are made available through Professor-mediated access, a direct comparison run is feasible. Second, the \texttt{OpenMed/OpenMed-NER-SpeciesDetect-BioMed-109M} model (DeBERTa-v3 + LoRA, SotA on Linnaeus and Species-800 benchmarks as of August 2025, arXiv 2508.01630) is now the recommended upgrade for organism NER without fine-tuning. A preprint describing a BioBERT fine-tune on 1,410 full-text microbiome articles (F1 = 0.96) is under peer review and will represent the best available public substitute once released (bioRxiv 2025.08.29.671515).

\subsection{Gap 2: Relation Extraction Model}

\textbf{MINERVA's approach:} BioMistral-AUG-7B—a 7-billion-parameter model fine-tuned on 1,100 labeled microbe-disease relation sentences, further augmented via Mixtral-8x7B paraphrase. Weights are not publicly released.

\textbf{This system's approach:} Gemini 2.5 Flash-Lite (zero-shot) via OpenAI-compatible API, using a 7-sample self-consistency protocol. Full agreement is required across all 7 samples before a relation is accepted. The measured full-consistency rate on the current corpus is \textbf{0.896}—meaning 89.6\% of relation candidates are confidently classified (accepted or rejected with unanimous agreement). This is a calibrated measure of extraction reliability, not a precision/recall comparison against labeled data.

\textbf{Residual risk:} Zero-shot LLM relation extraction performs poorly on specialized biomedical RE tasks when benchmarked against fine-tuned encoder models (Lever et al. 2023, BMC Bioinformatics). The self-consistency protocol mitigates this by requiring consensus, but does not eliminate domain-generalization error.

\textbf{Path to closure:} The highest-leverage upgrade is fine-tuning \texttt{microsoft/BiomedNLP-BiomedBERT-base-uncased-abstract} on the MicrobioRel dataset (bioRxiv 2025.08.03.666357), which covers 22 gut-microbiome relation types and yields F1 = 71.3\% for the best encoder model (arXiv 2506.08647). The Lever et al. 2023 labeled dataset (positive/negative/no association) is a second training option with fine-tuned BioLinkBERT achieving F1/P/R all above 0.80.

\subsection{Gap 3: UMLS Entity Normalization (and Beyond: Active TUI Gating)}

\textbf{MINERVA's approach:} Every entity is linked to a UMLS CUI via ScispaCy's EntityLinker. This enables alias deduplication—``NAFLD'' and ``nonalcoholic fatty liver disease'' resolve to the same node—and enables cross-graph merging using the UMLS ontology as a shared reference.

\textbf{This system's approach:} UMLS normalization is implemented and active. The \texttt{UMLSNormalizer} class uses ScispaCy's \texttt{scispacy\_linker} factory with the \texttt{en\_core\_sci\_lg} model (v0.5.4, spaCy 3.7.5) and the 2022 UMLS knowledge base (3,920,422 CUI entities). Both corpora have been enriched: the initial 640-paper corpus yields 96\% CUI coverage (52/54 disease entities, 54/54 microbe entities); the expanded 376-paper corpus yields 83\% coverage (115/131 disease, 125/131 microbe). Each relation pair in the graph carries CUI, TUI, similarity score, and official UMLS name fields. Disease entities are constrained to disease-class TUIs (\texttt{T019}, \texttt{T047}, \texttt{T191}, etc.); microbe entities to organism TUIs \{\texttt{T005} Virus, \texttt{T007} Bacterium, \texttt{T194} Archaeon, \texttt{T204} Eukaryote\}. Abbreviation resolution is enabled.

\textbf{Active TUI gating beyond MINERVA's spec.} This system goes one step further than MINERVA's documented use of UMLS: the same CUI/TUI metadata that MINERVA uses for alias deduplication is here promoted to an \emph{active acceptance gate} via the \texttt{EntityGate} module. A microbe surface form that fails to ground to a microbe-class TUI at similarity $\geq 0.85$ is rejected before relation classification, with the dropped record written to an audit JSONL with explicit \texttt{drop\_reason}. This catches entity-type errors that the upstream NER produced and that the substring vocab and self-consistency gate did not catch --- the canonical example being Edge \#5, where the gene-function noun phrase \texttt{gut bacterial ClpB-like gene function} was incorrectly tagged as an organism. A live audit of the eight currently-accepted text-derived microbe$\to$feature edges drops two: the Edge \#5 surface (similarity 0.799 grounded to ``Bacteria''~T007, below the 0.85 floor) and the misspelling \texttt{bacteriodetes} (no UMLS match). The 25\% drop rate sits below the 30\% recalibration threshold; the similarity floor is retained at 0.85.

\textbf{Residual risk:} UMLS similarity-based linking may produce false positives for short or ambiguous strings (e.g., ``CT'' linking to a drug CUI). The $\sim$17\% miss rate on the new-lanes corpus reflects entities whose surface forms are absent or poorly indexed in UMLS 2022 rather than a structural gap. Entities without a CUI match are retained in the graph as additive metadata; only the active \texttt{EntityGate} actually rejects, and only when the microbe-class TUI policy fails.

\textbf{Gap status: Closed and exceeded.} The architecture matches MINERVA's UMLS normalization approach for CUI-linked deduplication, and goes beyond it by activating the same UMLS metadata as a typed acceptance filter.

\subsection{Summary: Gap Status}

\begin{center}
\begin{tabular}{lll}
\textbf{Gap} & \textbf{Severity} & \textbf{Status} \\
\hline
NER model (microbe) & Medium & Best available substitute identified; \texttt{EntityGate} TUI \\
& & filter mitigates entity-type errors as a downstream stop \\
NER model (disease) & Low & BC5CDR baseline competitive; SotA upgrade identified \\
Relation model & Medium & Self-consistency rate 0.916; fine-tuning path documented \\
UMLS normalization & High & \textbf{Closed and exceeded} --- TUI-filtered, audit-active \\
Recall ceiling on text features & Medium & \textbf{Closed} --- BioClinical-ModernBERT dense retrieval \\
& & replaces 25-alias \texttt{\_FEATURE\_VOCAB} substring filter \\
Vision verifier monoculture & High & \textbf{Partially closed} --- dual gate alone is \\
& & structurally insufficient (both verifiers share \\
& & proposer bbox); 3 deterministic pre-verifier \\
& & gates added (caption / colorbar-detect / range \\
& & sanity with VLM colorbar tick extraction); \\
& & retroactive audit reduced 9 CORRELATES\_WITH \\
& & to 8 (1 vision + 7 text) \\
Measured precision/recall & High & In progress --- 66-row gold set sampled, \\
& & computer-aided Pass-1 labels generated; Pass-2 \\
& & after 14-day temporal window \\
\end{tabular}
\end{center}

\section{Limitations and Future Work}

\subsection{Radiomic Extraction Parameter Abstraction}

The current pipeline captures imaging modality and body location for each radiomic mention, but abstracts away low-level extraction parameters such as bin width, voxel resampling resolution, inter-slice spacing, and region-of-interest segmentation method. This is an inherent limitation of literature-mining relative to raw-data sharing: published papers rarely report these parameters in the abstract or structured body text that our pipeline processes. As a result, two GLCM entropy values extracted from the same organ on the same modality but with different bin widths are treated as equivalent nodes in the graph. Future work should incorporate parameter-aware feature identifiers, ideally by mapping extracted features to IBSI feature instance codes that encode extraction context, enabling downstream deduplication by both feature type and acquisition settings.

\subsection{Demographic Context and Population Heterogeneity}

Baseline body composition metrics—skeletal muscle index, visceral adipose tissue fraction, bone mineral density—are strongly influenced by age, sex, and BMI. The current graph aggregates evidence across heterogeneous clinical populations: the evidence supporting the edge ``\textit{Faecalibacterium} depletion associated with visceral adiposity'' may derive from studies in elderly patients post-hepatectomy and from studies in young obese adults, with potentially different phenotypic thresholds and effect sizes. Demographic stratification is not currently encoded as a node or edge property. Future iterations should add a \texttt{PopulationContext} property on edges or a \texttt{CohortNode} type to enable filtering by age range, sex, or BMI stratum, which would allow clinical users to identify microbiome-imaging phenotype associations that are consistent across demographic groups versus those that are cohort-specific.

\subsection{Inter-Scanner Variability and the Vision Track's Normalization Advantage}

A persistent challenge in multi-site radiomic studies is scanner variability: magnetic field strength, slice thickness, reconstruction kernel, and contrast agent protocol all affect absolute radiomic feature values and limit cross-study comparison of raw intensity measurements. The Vision Track's design addresses this at the evidence level. By extracting Pearson correlation coefficients ($r$-values) from heatmap figures rather than raw radiomic values, the Vision Track inherently normalizes across inter-scanner hardware variations. A correlation between a microbial taxon and GLCM Correlation, for instance, characterizes the \textit{rank-order relationship} between microbiome composition and texture structure---a property that is substantially more stable across scanners than an absolute HU value. This means that Vision Track edges carry a higher degree of cross-site generalizability than text-derived associations anchored to feature values from a single cohort.

\subsection{Interventions, Temporality, and Causal Graph Extension (Future Work)}

The current graph is observational: it encodes associations and correlations derived from cross-sectional or short follow-up studies, without distinguishing the direction of causality or capturing the effect of clinical interventions. An important next step is the integration of \texttt{Intervention} nodes---representing fecal microbiota transplantation (FMT), antibiotic courses, probiotic supplementation, or dietary modification---as a new node type connected to both microbial and imaging-phenotype nodes via typed \texttt{MODIFIES} or \texttt{REVERSES} edges. This would shift the graph from a purely observational network toward a partially causal model, enabling queries of the form: ``What imaging phenotypes are demonstrably reversible by FMT-mediated microbiome modulation?'' Achieving this would require longitudinal datasets with paired pre-/post-intervention microbiome and imaging assessments---a sparse but growing body of literature that should be targeted in future harvest lane design.

\subsection{Graph Sparsity and Evidence-Track Complementarity}

The delivered system produces three structurally distinct evidence layers that together close the intermediary gap. The \textbf{Text Track (phenotype $\to$ disease)} functions as a broad associative network: it mines phenotype co-mentions across 1,016 papers, yielding 74 \texttt{ASSOCIATED\_WITH} edges across 30 disease targets. The \textbf{Text Track (microbe $\to$ feature)} extracts direct associations between microbial entities and body composition features from the same entity-sentence corpus using Gemini self-consistency classification (7/7 full agreement required), producing 7 \texttt{CORRELATES\_WITH} edges. The \textbf{Vision Track} adds a strict quantitative backbone: a VLM proposes $r$-values from gradient heatmap figures, and a deterministic pixel-level verifier confirms them, producing 2 additional \texttt{CORRELATES\_WITH} edges with quantitative provenance.

Critically, the microbe $\to$ feature text extraction closes the three-hop path that the intermediary layer thesis requires. Three text-derived \texttt{CORRELATES\_WITH} edges connect to features that already carry downstream disease associations: \textit{Ruminococcus} $\to$ sarcopenia (37 diseases), \textit{Peptostreptococcus stomatis} $\to$ skeletal muscle index (7 diseases), and \textit{Eubacterium} $\to$ visceral adipose tissue (18 diseases). This yields \textbf{62 end-to-end traversable paths} of the form \texttt{Microbe $\to$ Feature $\to$ Disease}---the query pattern that MINERVA's single-hop \texttt{Microbe $\to$ Disease} schema cannot express. The remaining 6 \texttt{CORRELATES\_WITH} edges (4 body fat, 2 vision) currently lack downstream disease edges; these will close as the text track's phenotype extraction expands to additional feature vocabularies.

\end{document}
